\chapter{Semantic Search}
\label{ch:semantic}

Two questions come up the moment a store holds an ontology. \emph{What is under
this class?} is a question about structure, and the store can answer it
exactly. \emph{What is close in meaning to this term?} is a question about
similarity, and no index built so far in this book can answer it. The B+ trees
of \cref{ch:btree} are sorted, and it is tempting to think that sorted means
close. This chapter shows why that fails, what a hierarchy offers instead,
and where the vectors that do answer the second question live.

\section{What the tree orders}

A triple is a 24-byte key of three dictionary ids compared with
\code{memcmp} (\cref{ch:rdf-indexing}). The dictionary gives out ids densely, from
1, in the order terms are first seen: the memory dictionary takes
\code{next\_++}, and the disk dictionary takes \code{Size() + 1}. The loader
interns the subject, predicate, and object of each statement in file order. So
the order of every tree in OntoDB is \emph{first-seen order}. A class sits
next to another in the index only if their IRIs happened to appear close
together in the input. Nothing sorts by hierarchy, by label, or by meaning.

What the tree does give is \emph{structural} grouping: a bound prefix is one
contiguous range, which is why \code{?x rdf:type C} is a single scan of POS.

\section{A sort order cannot hold similarity}

Suppose we wanted one order on terms so that a range scan around a term returned
its neighbors in meaning. Embeddings live in hundreds of dimensions; the
question is already hard in two.

\begin{lemma}[Bandwidth]\label{lem:bandwidth}
Let $G$ be a graph on $n$ vertices with diameter $D$. Every linear order of the
vertices has an edge whose endpoints are at least $(n-1)/D$ apart in the order.
\end{lemma}
\begin{proof}
The first and the last vertex of the order are $n-1$ apart in it, and joined by
a path of at most $D$ edges. Some edge of that path covers at least $(n-1)/D$ of
the distance.
\end{proof}

Take the neighbor graph of a $k \times \dots \times k$ grid in $d$ dimensions:
$n = k^d$, diameter $d(k-1)$. Some two grid neighbors are at least
$(k^d-1)/(d(k-1)) \ge k^{d-1}/d$ positions apart, in \emph{any} order. A scan of
width $w$ around a point cannot hold all its neighbors unless $w$ is that large:
$\Omega(n^{(d-1)/d}/d)$. With a million points in two dimensions some neighbors
are about 500 positions away; in 384 dimensions the bound says ordering buys
nothing in the worst case. Space-filling curves (Z-order, Hilbert) keep
neighbors close on average and need several range scans to cover a ball; they do
not escape the lemma.

\begin{keyidea}
A B+ tree, or an adaptive radix tree, is a one-dimensional order. It is the
right structure for a prefix, a range of keys, or a subtree of a hierarchy. It is
the wrong structure for ``close in meaning''; that needs a vector index.
\end{keyidea}

\section{The exception: a hierarchy is a set of intervals}

A class hierarchy \emph{does} have an order in which everything below a class
is contiguous.

\begin{lemma}[Trees]\label{lem:tree}
Number the nodes of a rooted forest $1..n$ in postorder, and let $\mathit{lo}(c)$
be the smallest number in the subtree of $c$. Then $d$ lies in the subtree of $c$
if and only if $\mathit{lo}(c) \le \mathit{po}(d) \le \mathit{po}(c)$.
\end{lemma}
\begin{proof}
A postorder walk numbers a subtree as one contiguous block that ends at its
root.
\end{proof}

So ``all subclasses of $C$'' is one range scan over an index keyed by
$\mathit{po}$. Real ontologies are directed acyclic graphs: a class can have
several parents (the Gene Ontology has 57\,456 \code{is\_a} edges over 38\,092
terms). Following Agrawal, Borgida, and Jagadish \cite{agrawal1989}, choose a
spanning forest, number it in postorder, and give every class a \emph{set} of
intervals: its own tree interval merged with the intervals of every child.

\begin{lemma}[Interval sets]\label{lem:sets}
With $I(c) = \mathrm{merge}\big(\mathit{own}(c) \cup \bigcup_{c \to v} I(v)\big)$,
class $d$ is at or below $c$ if and only if $\mathit{po}(d)$ lies in an interval
of $I(c)$.
\end{lemma}
\begin{proof}
Induction over a reverse topological order: the numbers covered by $I(c)$ are
exactly those of the descendants of $c$.
\end{proof}

\begin{lemma}[Size]\label{lem:size}
Let $E'(c)$ be the non-tree edges whose tail is at or below $c$. Then
$|I(c)| \le 1 + |E'(c)|$.
\end{lemma}
\begin{proof}
A descendant $d$ is reached by a path from $c$. If the path uses no non-tree
edge, $d$ is in the tree subtree of $c$. Otherwise let $v$ be the head of the
last non-tree edge on the path; $d$ is in the tree subtree of $v$. The covered
set is therefore $\mathit{own}(c) \cup \bigcup_{e \in E'(c)} \mathit{own}(\mathrm{head}(e))$, a union of
$1 + |E'(c)|$ intervals, which has at most that many maximal intervals.
\end{proof}

A tree has no non-tree edges, so every class has one interval. Classes on a cycle
($A \sqsubseteq B \sqsubseteq A$) denote the same class; OntoDB collapses them
into one component before numbering. Which parent carries a class in the spanning
forest changes how many intervals the DAG needs; OntoDB takes the parent with the
longest path from a root.

\begin{figure}[ht]
\centering
\begin{tikzpicture}[font=\scriptsize,level distance=1.15cm,
  every node/.style={draw,rounded corners=1pt,inner sep=3pt,fill=white},
  level 1/.style={sibling distance=5.2cm},level 2/.style={sibling distance=2.4cm}]
\node (a) {A\quad$\mathit{po}=6$,\ [1,6]}
  child { node (b) {B\quad$3$,\ [1,3]} child { node (d) {D\quad$1$,\ [1]} } child { node (e) {E\quad$2$,\ [2]} } }
  child { node (c) {C\quad$5$,\ [2]\,[4,5]} child { node (f) {F\quad$4$,\ [4]} } };
\draw[-{Stealth[length=1.8mm]},dashed] (c.south west) -- node[draw=none,fill=none,midway,above,sloped]{non-tree edge} (e.north east);
\end{tikzpicture}
\caption{A DAG in which $E$ has two parents, $B$ and $C$. The spanning forest keeps
$E$ under $B$; the dashed edge is outside it. Each node shows its postorder number
and its interval set. $C$ needs two intervals: its own tree interval $[4,5]$ and
$[2]$, the label of $E$.}
\end{figure}

\subsection{Three ways to answer a type query}

Under RDFS, \code{?x a C} means ``$x$ has a type at or below $C$''.

\begin{center}
\begin{tabular}{@{}p{2.1cm}p{3.6cm}p{5.2cm}@{}}
\toprule
Strategy & Extra stored & One query costs \\
\midrule
Rewrite & nothing & one POS scan per class below $C$ \\
Materialize (\cref{ch:inference}) & a type triple per ancestor of each type & one POS scan \\
Intervals & one key per type triple & one range scan per interval of $C$ \\
\bottomrule
\end{tabular}
\end{center}

The interval index stores \code{(po(type), instance)} once per type triple. An
instance insert adds one key. The price is that labels change when the hierarchy
does, so OntoDB rebuilds them in $O(n+m)$ instead of maintaining them
(Agrawal et al.\ describe incremental maintenance with gaps; it is not built). A
related idea, numbering a tree in preorder and postorder so that axes become
regions of a plane, was used for XML by Grust \cite{grust2002}.

\begin{decision}{Ids stay first-seen; labels are a separate index}
\dpart{Choice} The dictionary keeps handing out dense ids in first-seen order. The
postorder labels live in their own B+ tree, \code{(label, member, 0)}, which
OntoDB rebuilds when the hierarchy changes.
\dpart{Alternatives} Renumber the dictionary so ids follow the hierarchy; then the
existing SPO, POS, and OSP trees would give subtree scans with no extra index.
\dpart{Why} Ids are stored in three permutations, in the log, and in snapshots.
A hierarchy edit would renumber all of them. A label is derived data.
\dpart{Cost} One more small tree, and a rebuild after the hierarchy changes.
\end{decision}

\section{Vector indexes}

The second question needs vectors. OntoDB keeps one unit-length float vector
per term id, so cosine similarity is a dot product. Embeddings come from
outside; a small sentence model on a CPU embeds the 38\,092 Gene Ontology names
in under a minute. OntoDB stores and searches them; it does not run the model.

\begin{description}[leftmargin=1.2em]
\item[Flat.] Score every vector: $O(nd)$ and exact. With 8-bit scalar
quantization (one scale per vector, $s = \max|x|/127$) a stored coordinate is
off by at most $s/2$, so for a unit query a score is off by at most
$(s/2)\sqrt d$, at a quarter of the memory.
\item[IVF] (inverted file \cite{jegou2011}). Cluster the vectors into $L$ lists
with $k$-means \cite{arthur2007}. A query scores the $L$ centroids, then only the
$p$ nearest lists: cost $Ld + p(n/L)d$, minimal at $L=\sqrt{pn}$, where it equals
$2d\sqrt{pn}$. OntoDB stores each list as a range of a B+ tree: key
\code{(list, row, 0)}, so a list is a prefix scan and an insert is a tree
insert.
\item[Product quantization.] Cut a vector into $m$ pieces and replace each by the
index of its nearest of 256 learned centers: $m$ bytes per vector. A query builds
$m$ tables of 256 dot products ($256d$ work), and a stored code scores as the sum
of $m$ table entries. That equals the dot product with the decoded vector exactly,
so all the error is in the stored side.
\item[HNSW] (hierarchical navigable small world \cite{malkov2020}). A layered
proximity graph searched greedily: memory $O(nM)$, the slowest to build, usually
the best recall for a given time, and no worst-case bound.
\end{description}

\begin{decision}{IVF lists are B+ tree ranges}
\dpart{Choice} A posting list is the set of keys \code{(list, row, 0)} in a
B+ tree of the same kind as the triple indexes.
\dpart{Alternatives} A separate array per list in memory, as FAISS keeps them
\cite{johnson2021}.
\dpart{Why} The tree already pages, caches, and latches; an insert or delete is a
tree operation; a list is a prefix scan, the access path OntoDB is built around.
\dpart{Cost} The vectors themselves are still fetched by row from the vector store,
one random access per candidate.
\end{decision}

\section{Filtering a vector search by a class}

A hybrid query asks for the $k$ terms nearest $t$ that lie under class $C$. Let
$s = |\mathrm{subtree}(C)|/n$.

\begin{description}[leftmargin=1.2em]
\item[Pre-filter.] Resolve $C$ to its members first (one range scan per
interval), then score only those, exactly: about $s\,n\,d$ work. Best when $s$ is
small.
\item[Post-filter.] Ask the index for more than $k$, keep those under $C$ (a
membership test is a binary search over $C$'s intervals; no scan), and ask for
four times as many until $k$ remain. If the filter were independent of the
similarity one would need about $k/s$ candidates, so it fails when $s$ is small.
\item[In-index filter.] Resolve $C$, then give the index the set as a mask: IVF
skips list entries outside it, HNSW still crosses masked nodes to stay connected
but never returns them \cite{patel2024}.
\end{description}

In an ontology the filter and the vectors are \emph{not} independent: the names
of a class's descendants are close to the class and to each other. A query from
inside $C$ finds most of its neighbors inside $C$, so post-filtering works far
better than $k/s$ predicts; a query from outside finds almost none, and it works
far worse. \Cref{sec:semantic-measure} measures both.

\section{Why not an ART with sibling links}

An adaptive radix tree \cite{leis2013} orders keys by byte, so its range scan is a
\emph{lexical} range scan: every label that starts with \code{rub}. That is a real
query, and the ART is good at it, but it is the same one-dimensional order as the
B+ tree, and \cref{lem:bandwidth} applies unchanged.

Sibling links on an ART's leaves are possible. Path compression and node growth
never move a leaf. Lazy expansion does: an insert that meets a leaf turns it into
an inner node and puts a \emph{copy} of the leaf under it, so a link that pointed
at the old object now points at an inner node. Every insert must also find its
sorted neighbors, a partial descent, and concurrent inserts must latch both. A B+
tree has the link for free: a leaf is a page, and a split is local. The
usual alternative is a stack-based iterator: no link, climb and descend a stack of
\code{[node, keys, next child]}, which pays pointer chasing between leaves and
costs nothing on insert. OntoDB's keys are three fixed-width
integers, which gives an ART little prefix to compress, so the B+ tree stays.

\paragraph{Pessimistic and optimistic prefixes.} Path compression stores a node's
shared prefix. \emph{Pessimistic} storage keeps the bytes in the node, so a
descent compares them and fails at the first wrong byte; the cost is
variable-length storage. \emph{Optimistic} storage keeps only the length: the
descent skips those bytes and the full key is checked once, at the leaf. Nodes are
smaller, and a wrong path is found late. For a range scan the difference is in the
seek. To find the first key at or above $\mathit{lo}$ the seek must know whether a
whole subtree lies below or above $\mathit{lo}$. With stored bytes it decides at the
node; with skipped bytes it must load a key from the subtree and compare. The
upper bound needs no jumping: a second seek finds the first key past $\mathit{hi}$.
The only binary search in a scan is over the sorted edge bytes of one node.

\section{What it costs}
\label{sec:semantic-measure}

These measurements are on the Gene Ontology (release 2026-07-26, without
obsolete terms: 38\,092 classes, 57\,456 \code{is\_a} edges), with the term names
embedded by a 22-million-parameter sentence model on a CPU. LUBM and a set of
synthetic instances give the type queries something to run on. Everything is
one thread on an Apple M2 Pro, and \code{benchmarks/semantic/run.sh} reproduces
it; the full tables are in \code{benchmarks/semantic/results.md}.

\paragraph{Does an order put neighbors together?} For 2\,000 random terms, the
ten nearest by cosine; \cref{tab:semantic-order} gives the share within $\pm W$
positions.

\begin{table}[ht]
\centering\small
\begin{tabular}{@{}lrrrrr@{}}
\toprule
Order & median gap & $\pm$100 & $\pm$1000 & $\pm$5000 \\
\midrule
dictionary id (first seen) & 1739 & .211 & .284 & .438 \\
hierarchy postorder & 45 & .377 & .561 & .735 \\
label, alphabetical & 1749 & .334 & .408 & .558 \\
random & 10823 & .005 & .049 & .242 \\
\bottomrule
\end{tabular}
\caption{Share of the ten nearest neighbors found within $\pm W$ positions.}
\label{tab:semantic-order}
\end{table}

The postorder of the hierarchy is the best order, and far better than first-seen
ids: the nearest neighbor is a median of 45 positions away instead of 1\,739. It
still finds only 38\% of the ten within 100 positions, and 56\% within 1\,000, as
\cref{lem:bandwidth} says to expect. First-seen ids are not random either: the
Gene Ontology minted its ids in related batches, so id order carries locality that
belongs to the file, not to the engine. A term's neighbor is a parent or child in
12\% of cases and within two \code{is\_a} steps in 32\%, in a two-step
neighborhood of 45 terms (0.12\% of all terms): the hierarchy is informative about
meaning without being sufficient.

\paragraph{Intervals.} The Gene Ontology needs 76\,570 intervals for 38\,092
classes: 2.01 per class on average, 598~KiB in all, at most 741 for one class,
inside the bound of \cref{lem:size} (19\,368). Materializing the subclass closure
would store 393\,448 ancestor pairs. Whether $d$ is below $c$ takes 27~ns, 29 times
faster than a search up the parents. The descendants of a class with 1\,000 or more
subclasses take 101 range scans and 40 page reads, against 3\,304 scans and 694
pages for one scan per class.

\paragraph{Three ways to answer \code{?x a C}.} On 200\,000 synthetic entities over
the Gene Ontology's classes (495\,552 triples):

\begin{center}
\small
\begin{tabular}{@{}lrrr@{}}
\toprule
Strategy & triples & file (MiB) & load (s) \\
\midrule
rewrite & 495\,552 & 63.6 & 5.1 \\
intervals & 495\,552 & 72.9 & 5.3 \\
materialize & 4\,017\,973 & 309.5 & 180.3 \\
\bottomrule
\end{tabular}
\end{center}

Materializing stores 8.1 times the triples and takes 36 times as long to load. For
a class with 1\,000 or more subclasses a query reads 265 pages with intervals, 182
with the materialized store, and 12\,803 by rewriting (48 times more than
intervals). All three return the same answers on every class sampled. On LUBM(1),
whose hierarchy has 43 classes and little depth, the three are within a factor of
about two; intervals have little to save where the hierarchy is shallow.

\begin{decision}{Intervals instead of a materialized closure for type queries}
\dpart{Choice} A query for the instances of a class is one range scan per interval
of the class, over an index of \code{(label of type, instance)}.
\dpart{Alternatives} Materialize the entailed type triples (\cref{ch:inference});
rewrite into one scan per subclass.
\dpart{Why} On a deep, wide hierarchy the interval index costs 15\% more space
instead of 4.9 times as much, loads in 5 seconds instead of 180, and reads within 1.8
times of the materialized store's pages.
\dpart{Cost} The labels are rebuilt when the hierarchy changes, and only
\code{subClassOf} is understood: no \code{subPropertyOf}, \code{domain}, or
\code{range}. A materialized closure handles those rules with the same machinery
and an interval index does not.
\end{decision}

\paragraph{Vector indexes.} 38\,092 vectors of 384 floats; recall@10 against the exact
answer, time per query on one thread.

\begin{center}
\small
\begin{tabular}{@{}llrrr@{}}
\toprule
Index & setting & recall & $\mu$s & MiB \\
\midrule
flat float32 & exact & 1.000 & 2\,238 & 55.8 \\
flat PQ $m{=}48$ & codes only & 0.778 & 819 & 1.7 \\
IVF $L{=}512$ & nprobe 16 & 0.929 & 255 & 56.5 \\
IVF $L{=}512$ & nprobe 64 & 0.976 & 928 & 56.5 \\
HNSW $M{=}16$ & $ef$ 10 & 0.993 & 44 & 59.7 \\
HNSW $M{=}16$ & $ef$ 40 & 0.997 & 108 & 59.7 \\
\bottomrule
\end{tabular}
\end{center}

HNSW dominates on this data: 0.993 recall at 44~$\mu$s, 51 times faster than the
exact scan. A word on that baseline: the first run scored with a plain dot-product
loop, which a compiler may not vectorize without \code{-ffast-math}, and the exact
scan took 8.3~ms. With eight independent accumulators it takes 2.2~ms, and a
speed-up over ``exact'' is only as meaningful as the exact baseline is fast. IVF
lists stored in the B+ tree cost about 2.7 times as much per scored row as a flat
scan (157~ns against 59~ns): a decoded key and a random read of the vector, the
price of lists that page and update like the triples. Product quantization is for
memory; re-ranking its best 100 exactly gives 0.998 recall, but needs the floats.

\paragraph{The ten nearest terms under a class.} \Cref{tab:semantic-hybrid} gives
time and recall for five subtree sizes, with a query from inside the class.

\begin{table}[ht]
\centering\small
\begin{tabular}{@{}lrrrrr@{}}
\toprule
Members of $C$ & $\le$10 & 11--100 & 101--1000 & 1001--5000 & $>$5000 \\
\midrule
pre-filter & 3 & 9 & 62 & 292 & 1\,404 \\
post-filter, HNSW & 37\,553 & 1\,739 & 427 & 70 & 68 \\
masked HNSW & 16\,653 & 15\,846 & 1\,150 & 183 & 105 \\
masked IVF & 1\,253 & 111 & 104 & 124 & 161 \\
auto & 3 & 9 & 59 & 269 & 95 \\
\bottomrule
\end{tabular}
\caption{Microseconds per query, query inside $C$. Pre-filter, post-filtered and masked HNSW are exact at these sizes; masked IVF has recall 0.95 to 1.00.}
\label{tab:semantic-hybrid}
\end{table}

Pre-filtering through the interval index is exact and the fastest up to about a
thousand members, by factors from 1.7 to more than 10\,000. The filter is
correlated with the vectors: of the unfiltered top 100 neighbors of a term inside
$C$, 4\% to 72\% lie under $C$, and for a term outside it 0\% to 12\%. The
independence estimate $k/s$ is therefore wrong in both directions, and
post-filtering IVF recalls 0.07 of the right answers on the smallest classes for a
query outside them. A masked HNSW search wins above 5\,000 members (105~$\mu$s inside
and 900~$\mu$s outside, exact). One more lesson: resolving the members into a bitmap
before a masked search cost 560~$\mu$s of a 669~$\mu$s query, while testing each
touched row's label against the intervals, one binary search, cost nothing extra
and brought it to 105~$\mu$s. \code{kAuto} takes the exact path up to 4\,096 members
and a masked HNSW search above, and is within a factor of about two of the best
strategy everywhere but one cell.

\begin{decision}{Resolve small classes, mask large ones}
\dpart{Choice} Up to 4\,096 members: resolve the members by range scans and score
them exactly. Above: search HNSW with a predicate, ``the row's label lies in an
interval of $C$''.
\dpart{Alternatives} Always post-filter; always mask; always resolve.
\dpart{Why} Resolving is exact and fastest while a class is small. Post-filtering
fails for a query outside a small class. The threshold sits below the point where an
outside query makes masking slower than resolving (7\,000 to 25\,000 members) and
loses at most a factor of 3.8 where an inside query makes masking faster (about
1\,000).
\dpart{Cost} One threshold for both kinds of query; a class of 1\,001 to 5\,000
members queried from inside runs 3.8 times slower than it could.
\end{decision}


\begin{status}
Built (plan 14.1--14.9): class labels, the interval indexes, the vector store,
flat, IVF, and HNSW indexes, hybrid search, the shell commands
(\code{\textbackslash hindex}, \code{\textbackslash under},
\code{\textbackslash vload}, \code{\textbackslash vindex},
\code{\textbackslash similar}), and the benchmarks of this chapter
(\code{benchmarks/semantic/}). Vectors live in memory, not in buffer pool pages,
and the structures are not updated by inserts: they are rebuilt.
\end{status}

\begin{checkyourself}
\item Why does a sort order of dictionary ids not put similar classes next to each other? What would have to be true of the input file for it to?
\item Use \cref{lem:bandwidth} to bound how far apart two neighbors must be in the best order of a $10 \times 10$ grid.
\item Give the postorder labels and the interval sets of the diamond $D \sqsubseteq B, C \sqsubseteq A$ for a spanning forest that keeps $D$ under $B$.
\item A class has 5 non-tree edges below it. How many intervals can it have at most, and why?
\item When is pre-filtering better than post-filtering for a class filter, and when is it not?
\item Why can post-filtering work better than the independence estimate $k/s$ says, in an ontology?
\item What does a seek for a lower bound have to do differently if prefixes are stored optimistically?
\end{checkyourself}

\section*{Further reading}
Agrawal, Borgida, and Jagadish \cite{agrawal1989} on interval labelling of
transitive relationships; Jégou, Douze, and Schmid \cite{jegou2011} on
inverted files and product quantization; Malkov and Yashunin \cite{malkov2020}
on HNSW; Patel et al.\ \cite{patel2024} on filtered vector search; Pan, Wang, and
Li \cite{pan2024} for a survey of vector database systems; Leis, Kemper, and
Neumann \cite{leis2013} on the adaptive radix tree.
